\subsection{Measured Response: Predictive Scope of Positional Profiles}\label{sec:spectralpredictionresults}

Both frozen models reproduce the observed signs of the nonce-word gaps at all six cipher-by-width settings, including the small positive 128-bit ChaCha20 gap, for which no sign was required. Figure~\ref{fig:mechanisticextensions}a,b compares the frozen predictions with the new observations. Agreement in direction does not ensure agreement in magnitude or spectrum. At 128 bits, the observed ChaCha20 fundamental contrast is $0.53850$: M0 predicts $0.42303$, whereas M1 predicts $0.55252$. The corresponding Salsa20 values are $0.80082$, $0.76860$, and $0.80452$. Retaining conditional training-key and base profiles therefore moves both spectral predictions closer to these observations, while leaving an excess fundamental share in M1.

Across all widths, the mean AC-spectrum $L^1$ errors decrease from $0.322$--$0.379$ under M0 to $0.078$--$0.084$ under M1 for ChaCha20, and from $0.101$--$0.261$ to $0.010$--$0.018$ for Salsa20. M1 also has smaller absolute fundamental-share errors in all six settings. Its advantage does not extend to every scalar summary: M0 has smaller absolute $G$ errors for all three ChaCha20 widths. For Salsa20, M1 reduces the gap error at each width; at 128 bits, the observed $G=0.80132$ bits compares with predictions $0.85574$ and $0.79465$ from M0 and M1.

The archived tables report signed errors, conditional whole-key intervals and Monte Carlo standard errors for both models. These empirical models separate the contribution of retained conditional profiles within the specified construction. Remaining discrepancies can include finite training-fit error and omitted residual dependence; the comparison does not identify their individual contributions. Errors are reported quantitatively without a model-rejection test, and small average spectral errors do not guarantee accurate predictions for individual masks or keys.
